\section{The relative picture, and why it is the right one}\label{sec:relative}

Everything computed so far lives in the absolute metric of
\cite{graf2026numerical,noferini2016numerical}, and deliberately so: it is the metric in
which the Gaussian calculus is exact, every average is an identity or a quadrature, and
the machinery of Sections~\ref{sec:jacobian}--\ref{sec:structural} applies without
normalizing anything. But it is not the metric in which conclusions should be drawn.
Floating-point arithmetic commits \emph{relative} errors; backward-stable eigensolvers
return backward errors relative to the norm of the data; and the relative condition
number is invariant under the rescalings of data and unknowns that every implementation
performs, while the absolute one is not. In Gaussian ensembles this is not a pedantic
distinction: the hidden scales are enormous and root-dependent, as the outer-root and
units effects documented throughout this paper show, so absolute comparisons
systematically conflate genuine sensitivity with the bookkeeping of scales. The absolute
metric is the right vehicle for the computation; the relative metric is the right lens
for the verdict. This section converts the entire scoreboard.

Following \cite[Sec.~5]{noferini2016numerical},
$\kappa_{\mathrm{rel}} = \kappa_{\mathrm{abs}}\cdot\|D\|/\|X\|$ with $D$ the data of the
subproblem and $X$ its solution, so for every method
\[
  \what\Rcal_M := \frac1N\,\E\sum_j\log
  \frac{\kappa_{\mathrm{rel}}(M)}{\kappa_{\mathrm{rel,root}}}
  = \Rcal_M + \E\log\|D_M\| - \E\log\max_i\sup_\Omega|p_i| + O(1),
\]
the output scales contributing $O(1)$ in these models, and the baseline being
$\E\log\max_i\sup_\Omega|p_i| = \frac d2\log n + O(\log dn)$.

\begin{proposition}[relative scoreboard]\label{prop:scoreboard}
  In the separate-degree model, with the pairwise and local entries at $d=2$:
  \begin{enumerate}[(a)]
    \item Cayley: $\what\Rcal = O_d(\log\log n)$ for every $d$, by
      Proposition~\ref{prop:opnorm}: relative neutrality, the clean case.
    \item Sylvester: $\what\Rcal = \frac12\log n + O(\log\log n)$: the data scale
      $\sup_x\|S(x)\|_2$ coincides with $\sup|p|$ up to the Toeplitz-symbol log factors,
      so the relative and absolute coefficients agree.
    \item Multiparameter (\textsc{MinUnif}): $\what\Rcal = \frac{13}2\log n + O(1)$, the
      data norm contributing $\frac12\log n$ through the coefficient block.
    \item Normal form:
      $\what\Rcal = (\frac d2+1)\log n + \E\log\|M_g(\B)\| + O(1)$, basis-dependent
      through the norm of the multiplication matrix; measured,
      $\E\log\|M_g\|\approx\frac12\log N$ for the companion matrix at $d=1$ and
      $\approx2\log n$ for the B\'ezout pencil at $d=2$, giving $\approx2\log n$ and
      $\approx4\log n$ respectively.
    \item Gr\"obner elimination: $\what\Rcal = -n(\log n+C)(1+o(1))$: only \emph{half}
      of the pair energy cancels against $\E\log\|f\|$, because the coefficient norm of
      the eliminant sees one radial tail, $\sum_i\E\log\max(1,|x_i|)$, where the
      root-averaged potential sees two; measured,
      $\log\|f\|/\mathrm{PE} = 0.65, 0.64, 0.62$ at $n=6,8,11$, trending to $\frac12$.
    \item Rational univariate representation: $\what\Rcal = \delta\,N\,(1+o(1))$ with,
      in the balanced case,
      \[
        \delta = \frac G\pi + \frac{\log2}4 - \frac{7\,\zeta(3)}{2\pi^2}
        = 0.03857\ldots > 0,
      \]
      $G$ Catalan's constant: the unique exponential penalty of the family.
  \end{enumerate}
\end{proposition}

\proofsketch
(a) is Proposition~\ref{prop:opnorm} with the conversion formula. (b), (c), (d) insert
the data scales computed or measured above; for (d) note that $\|M_g(\B)\|$ is a
property of the chosen basis, not of the quotient. For (e), $\E\log\|f\|$ per root is
the circle average of the potential of the coordinate marginal,
$\E\log\max(1,|x|) = \frac1{2n}(\log n+C')(1+o(1))$, a single-tail quantity, while the
pair energy is the root average of the same potential, a two-tail quantity twice as
large; the difference is $-n(\log n+C)$ at leading order. For (f), the same comparison
for the projected marginal is a contest of constants: the circle potential of the limit
law $W$ is
$c' = \E[\log_+|W|] = \frac1{2\pi}\int_{-\pi/2}^{\pi/2}\log(1+\cos\theta)\,d\theta
= \frac G\pi - \frac{\log2}4$, by the classical evaluation
$\int_0^{\pi/2}\log(1+\cos\theta)\,d\theta = 2G - \frac\pi2\log2$, while the pair energy
constant is $I(u) = \frac{7\zeta(3)}{2\pi^2} - \frac{\log2}2$
(Proposition~\ref{prop:rur}); $\delta = c' - I(u)$. Measured,
$(\log\|f\| - \mathrm{PE})/N = -0.073, -0.052, -0.029$ at $n=6,8,11$, rising toward
$\delta$ with the explicit $(\log n)/n$ outer transient common to both terms.
\endproof

The relative verdict is sharper than the absolute one, and closer to practical folklore.
The resultant methods are benign: Cayley exactly neutral, Sylvester polynomially
penalized, and the exponential worst cases of \cite{noferini2016numerical} are tail
events in both metrics. The normal form pays a polynomial factor that depends on the
basis through $\|M_g\|$, one more reason the basis is not a detail. Elimination remains
distorted by units even relatively, now through the mismatch between the unit circle,
where the coefficient norm of the eliminant is measured, and the region where its roots
actually live. And the rational univariate representation is the one method the relative
metric convicts: its projected root measure fails to be an equilibrium measure of the
circle, its potential is not constant, and the resulting $e^{\delta N}$ average penalty
is exactly the clustering mechanism of the devastating example of
\cite[Sec.~4]{graf2026numerical} surviving into the average case, in the metric where it
belongs. The practical reading: project only if you must, and if you must, the constant
$\delta(u)$ is a design objective, minimized by making the projected measure as close to
equilibrium as the geometry allows.
